\subsection{問題分析}
\subsubsection{車輛投放問題}
\begin{enumerate}
    \item 投放選址\\
    車輛投放選址間須有合適的距離，否則會導致不必要的投放位置，同時，像是機場這樣的不可投放區域也須考慮在內.
    \item 初始投放數量\\
    車輛若未根據區域特徵（如居民區、旅遊區、樞紐區）的人口密度、出行需求精準投放，會導致部分區域“車滿為患”造成車輛堆積、佔用公共空間，而另一些區域則“一車難求”的問題.
\end{enumerate}
車輛投放最佳化： 解決「在哪裡放多少車」的問題. 需基於歷史騎行大數據、城市POI資訊、人口分佈、土地利用、公共交通網絡等，建構需求預測模型，科學規劃不同區域、不同時間點的初始車輛投放配額和停車點容量設置，力求從源頭減少嚴重不均衡.
  
\subsubsection{車輛調度問題}

\subsubsection{共享單車的再平衡問題}
本問題核心在於優化共享單車運營商的卡車調度策略，以最小化總成本（運輸成本 + 懲罰成本）實現站點間單車的再平衡。其本質是一個帶有懲罰機制、容量約束和時間窗限制的取送貨車輛路徑問題。解決問題的基本思想與關鍵點分析如下：

\begin{enumerate}
    \item \textbf{核心目標與權衡：}
    \begin{itemize}
        \item \textbf{根本目標：} 最小化總運營成本
        \item \textbf{核心矛盾：} 完全滿足所有站點需求可能需要極高的運輸成本（長距離、多卡車、超時工作），甚至因物理限制（卡車容量、時間）無法實現
        \item \textbf{關鍵權衡：} 模型允許進行「成本-服務質量」權衡：
        \begin{itemize}
            \item 支付運輸成本：派遣卡車運輸單車，直接調整站點庫存使其接近期望需求
            \item 支付懲罰成本：不派遣卡車訪問某些站點或裝卸不足，容忍庫存偏差，承擔用戶不便成本
        \end{itemize}
        \item \textbf{建模思想：} 通過聯合優化卡車的路線選擇、裝卸量決策以及可接受的庫存偏差，在滿足物理和運營約束前提下，找到運輸成本與懲罰成本之和最小的方案
    \end{itemize}

    \item \textbf{問題分解與建模框架：}
    \begin{itemize}
        \item \textbf{基礎模型（單卡車）：} 解決基本情形——僅用一輛卡車調度：
        \begin{itemize}
            \item \textbf{路徑：} 卡車從車場出發，訪問站點（每站最多訪問一次）後返回車場
            \item \textbf{裝卸操作：}
            \begin{itemize}
                \item 取貨點：取車量不超過站點初始庫存
                \item 送貨點：卸車量不超過站點空位容量
            \end{itemize}
            \item \textbf{卡車載量：} 確保載量連續性與容量限制（不超過卡車最大容量）
            \item \textbf{時間約束：} 累計工作時間（行駛時間+裝卸時間）不超過預定工作時間
            \item \textbf{庫存與懲罰：} 最終庫存偏差帶來懲罰成本
        \end{itemize}
        \item \textbf{擴展模型（多卡車）：}
        \begin{itemize}
            \item 為每輛卡車複製單卡車模型的核心決策變量和約束條件
            \item \textbf{耦合因素：}
            \begin{itemize}
                \item 站點最終庫存：所有卡車在該站點操作的總效果
                \item 車場資源：假設車場有無限單車（初始載量無上限約束）
            \end{itemize}
            \item \textbf{目標函數：} 最小化所有卡車運輸成本之和加上所有站點懲罰成本之和
        \end{itemize}
    \end{itemize}

    \item \textbf{關鍵約束的處理思想：}
    \begin{itemize}
        \item \textbf{大數值法：} 處理邏輯條件（如「若選擇某路段則需滿足載量約束，否則該約束失效」）
        \item \textbf{站點分類與懲罰權重：} 通過差異化懲罰成本（樞紐區=0，居民區較低，旅遊區較高），引導模型優先滿足對庫存偏差容忍度低的區域
    \end{itemize}

    \item \textbf{參數獲取與實踐意義：}
    \begin{itemize}
        \item \textbf{數據驅動：} 參數（站點位置/類型/容量/庫存需求、道路距離/時間、卡車屬性、懲罰權重、裝卸效率）需基於實際運營數據和地圖服務獲取
        \item \textbf{懲罰成本的核心作用：} 連接運營決策（調度成本）與用戶體驗（需求未滿足帶來的不便）
        \item \textbf{時間與容量約束：} 直接反映調度的物理可行性和運營成本
    \end{itemize}
\end{enumerate}

\textbf{基本思想：} 將共享單車再平衡問題抽象為帶懲罰的多約束車輛路徑優化問題。通過引入可接受的庫存偏差及對應的懲罰成本，在滿足卡車物理限制和站點操作約束前提下，系統性權衡直接運輸成本與服務不完美導致的間接成本。利用圖論建模站點和路徑，通過整數規劃和線性規劃技術描述決策變量，並通過站點分類懲罰機制實現區域服務優先級差異化處理。單卡車模型是基礎構件，多卡車模型通過複製和耦合全局懲罰項實現擴展。最終目標是求得總成本最低的卡車調度方案（路線+裝卸量）。